% Figure: discontinuity (edge-bending) stress at a cylinder-to-head
% junction. The membrane solution is uniform; the local bending boundary
% layer adds a stress that decays over a length scale sqrt(R h). Plot of
% the total meridional stress against axial distance from the junction,
% normalised by the membrane value.
\begin{figure}[htbp]
\centering
\begin{tikzpicture}
\begin{axis}[
  width=11cm, height=6.5cm,
  xlabel={axial distance from junction $\ x/\sqrt{R h}$},
  ylabel={meridional stress / membrane stress},
  xmin=0, xmax=6, ymin=0, ymax=2.0,
  axis lines=left,
  grid=major, grid style={enggray!20},
  tick label style={font=\scriptsize},
  label style={font=\small},
  legend style={font=\scriptsize}, legend pos=north east,
]
% membrane reference
\addplot[enggray, dashed, very thick, domain=0:6, samples=2] {1.0};
\addlegendentry{membrane only}
% total stress: membrane + decaying bending boundary layer.
% schematic damped form 1 + 0.8*exp(-x)*cos(x), the classical edge
% bending decay e^{-beta x}(cos + ...).
\addplot[engblue, very thick, domain=0:6, samples=121]
  {1 + 0.8*exp(-x)*cos(deg(x))};
\addlegendentry{membrane $+$ bending}
% mark the peak near the junction
\node[engred, font=\scriptsize, anchor=south west] at (axis cs:0.15,1.8)
  {edge peak};
\draw[engred, ->] (axis cs:0.5,1.78) -- (axis cs:0.05,1.78);
% mark the decay length
\node[enggray, font=\scriptsize, anchor=south] at (axis cs:3.2,1.06)
  {decays over $\sim \sqrt{R h}$};
\end{axis}
\end{tikzpicture}
\caption[Discontinuity stress at a shell junction]{Meridional stress
near a cylinder-to-head junction, normalised by the membrane stress
(schematic). Membrane theory predicts the uniform dashed line; the
real junction must also satisfy continuity of radial displacement and
slope between two shells of different stiffness, which it does by a
local bending boundary layer (solid curve). The bending stress peaks
at the junction, can reach a sizeable fraction of the membrane stress,
and decays over an axial length of order $\sqrt{R h}$. The peak is why
a knuckle radius, a thickness taper, or both are detailed at every
head-to-shell junction, and why a pure membrane sizing is necessary
but not sufficient.}
\label{fig:vol03ch09:discontinuity-stress}
\end{figure}
